\section{A rational certificate for the degree-31 example}
\label{cert:section}

This appendix proves the root bounds in Proposition~\ref{np:roots} and
the predictor estimates in Proposition~\ref{np:moments}.  The root
calculation uses exact rational arithmetic.  The moment calculation
combines finite rational products with an analytic estimate for their
infinite tails.  We specify the rounding and the order of operations so
that the rational weights in~\eqref{cert:weights} are completely determined.
The accompanying script \texttt{verification/verify\_degree31.py} implements
these instructions using only Python integers and rational fractions;
the recurrences below give the complete certificate independently of
that implementation.  Every terminating decimal in this appendix
denotes its exact rational value.

\subsection{Locating all the roots}
\label{cert:roots}\label{cert:root-proof}

Recall the polynomial
\[
 p(t)=\sum_{j=0}^{9}(-t^{3j}+t^{3j+1})-2t^{30}+t^{31}.
\]
Put $E=10^{40}$.  The entries in Table~\ref{cert:root-table} are the real
and imaginary numerators of centers $z=(a+\mathrm i b)/E$.
Include the conjugate of each nonreal center.  Thus the table specifies
$31$ disks, all of radius $\delta=10^{-35}$.
Denote its first and sixth centers by $z_\beta$ and $z_\alpha$,
respectively.

\begin{table}[htbp]
\centering
\begingroup\scriptsize\ttfamily
\begin{tabular}{rr}
\normalfont Real numerator & \normalfont Imaginary numerator\\\hline
18392867573155250211326543340571993487329 & 0\\
-9741384652837567528189111907762645564212 & 960611633124778624015947511546714789832\\
-9390979934023233037686839914141181659013 & 2848063934279278841200845525848576791934\\
-8707010393241200595697704637503228167256 & 4634358961359777463662031763152574795910\\
-7723118558199735333585947045021501410451 & 6250214494152701297733561366714273855663\\
-6480873022754265587670654805348739312397 & 7617033927823246411701558316834470378693\\
-3296776282962928450195355562079974291919 & 9362810613239308448752821856735617305389\\
-1442513481645415489053007250852371013818 & 9733039088635544125591234293715446465839\\
473579005983151143865690217073943282984 & 9759329295153416086962240697595385836176\\
2372858204750359429049351486827380431084 & 9430908396110323024329804770889268443510\\
4184061662454447896993374181551299054853 & 8750376307352230997951514116951420713869\\
5840543823426662247726359012202878127980 & 7733999914049269380846159864123679516166\\
7280042164772721585016920411987150804420 & 6409668313572020921539025901166241650504\\
8444162922855313595050499086097221186662 & 4814843758334683329011653039104490553475\\
9275649862302447872891854755866108397375 & 2996533501299661135726051235425500195711\\
9715324892541617145821300300817663390044 & 1019107767839771580294306582422817927060
\end{tabular}
\endgroup
\caption{Rational centers for the roots; both columns have denominator $10^{40}$.}
\label{cert:root-table}
\end{table}

For a complex number write $\|z\|_1=|\Re z|+|\Im z|$ and
$\|z\|_\infty=\max(|\Re z|,|\Im z|)$.  Exact Horner evaluation gives,
at every displayed center,
\begin{equation}\label{cert:residual}
 |z|<2,\qquad \|p'(z)\|_\infty>1,\qquad
 2\cdot10^{35}\|p(z)\|_1<\|p'(z)\|_\infty.
\end{equation}
Here is an integer prescription for these comparisons.  For
$q(t)=\sum_{j=0}^d q_jt^j$ and the Gaussian integer $Z=Ez$, initialize
$(V,h)=(q_d,1)$ and, for $j=d-1,d-2,\ldots,0$, replace
\[
 h\leftarrow Eh,\qquad V\leftarrow ZV+hq_j.
\]
The final value is $q(z)=V/h$.  Apply this to $p$ and $p'$, and clear
the positive denominators in~\eqref{cert:residual}.  Comparisons of
squared distances give $|z-z'|>.01$ for distinct centers, including
conjugates.  Comparisons of coordinates and squared moduli also give
\begin{equation}\label{cert:center-bounds}
 \begin{gathered}
 1.83928<z_\beta<1.83929,\qquad
 1.0002090<|z_\alpha|^2<1.0002093,\\
 |z|<.9939\quad\hbox{at all other displayed centers}.
 \end{gathered}
\end{equation}
For $|h|=\delta$, the terms of degree at least two in
$p(z+h)-p(z)-p'(z)h$ have total modulus at most
\[
 \delta^2\sum_{j=0}^{31}|p_j|3^j<10^{17}\delta^2.
\]
Indeed, $|z|<2$ and $\delta<1$, so the binomial expansion gives this
bound term by term.  By~\eqref{cert:residual}, the constant term is
less than $|p'(z)|\delta/2$, whereas the displayed remainder is less
than $|p'(z)|\delta/2$.  Rouch\'e's theorem therefore puts exactly one
root, counted with multiplicity, in each disk.  The disks are disjoint
and their number is the degree, so these are all the roots and all are
simple.  The root in the real-centered disk is real by conjugation.
Every other disk misses the real axis, so the remaining roots form
fifteen nonreal conjugate pairs.
The radius $10^{-35}$ and~\eqref{cert:center-bounds} imply every bound
in Proposition~\ref{np:roots}, including $|1/\alpha|<.9999$.

\subsection{Rounded coefficient arrays and finite products}
\label{cert:algorithm}

The numerical inputs are the sequences $C_0,C_1$ defined in the
main text in~\eqref{np:coefficients}.  We approximate their single and paired moments
\begin{equation}\label{cert:exact-moments}
 P_A=\prod_{m\in\mathbb Z}\cos(\pi C_A(m)),\qquad
 P_{AB}^{\sigma}(k)=
 \prod_{m\in\mathbb Z}\cos\bigl(\pi(C_A(m)+\sigma C_B(m+k))\bigr),
\end{equation}
where $A,B\in\{0,1\}$, $\sigma\in\{1,-1\}$, and $k\ge0$.
Both products converge by exponential decay of the coefficients.
The sign notation $+$ and $-$ means $\sigma=1$ and $\sigma=-1$.

We now specify the approximations used to define the weights.  Use
\[
 S=10^{30},\qquad R=3000,\qquad H=100,\qquad K=2000,
 \qquad W=10^9.
\]
For a real rational $x$ set
\[
 [x]=\frac{\lfloor Sx+1/2\rfloor}{S},
\]
and round complex numbers componentwise.  Thus, if $x=u/v$ with
integers $u,v$ and $v>0$, its rounded numerator is
$\lfloor(2Su+v)/(2v)\rfloor$.  Additions, subtractions, conjugation,
and multiplication by an integer are exact unless brackets are
explicitly displayed.  Set
\[
 \pi'=\frac{3141592653589793238462643383279}{S}.
\]
The inequality $|\pi'-\pi|<1/S$ follows by using $30$ terms of the
alternating series for $\arctan(1/5)$ and $10$ for $\arctan(1/239)$
in Machin's identity
$\pi=16\arctan(1/5)-4\arctan(1/239)$.  The total remainder bound is
\[
 \frac{16}{61\cdot5^{61}}+\frac4{21\cdot239^{21}}.
\]
In particular this verification, too, uses rational arithmetic only.

The contracting and expanding roots are denoted by $I$ and $D$, as in
the main text.  For each representative center $z$ first form
$r_z=[1/p'(z)]$, with $p'(z)$ evaluated exactly.  Store a base and a
coefficient as follows:
\begin{equation}\label{cert:base-pairs}
 (a_z,v_z)=
 \begin{cases}
 ([z],r_z),&|z|<1,\\
 ([1/z],-[[1/z]r_z]),&|z|>1.
 \end{cases}
\end{equation}
For a contracting representative, start $v=v_z$ and add $2\Re v$
to $c_0(-1)$, then to $c_0(-2),c_0(-3),\ldots$, replacing
$v\leftarrow[va_z]$ after each addition.  For an expanding
representative do the same on $c_0(0),c_0(1),\ldots$; the multiplier
is $1$ for the real root and $2$ for the nonreal representative.
All arrays start at zero.  Compute $c_0(m)$ on
\[
 -R-K-5\le m\le H+K+5
\]
and define, wherever these three entries are present,
\begin{equation}\label{cert:shift-array}
 c_1(m)=c_0(m+3)-c_0(m)+c_0(m-1).
\end{equation}
These extra endpoint entries ensure that every subsequent use of
$c_1$ lies in its defined range.  The arrays approximate the sequences
$C_0,C_1$ in the main text.

For a rational input $c$ let $g(c),h(c)$ be the following rounded
cosine and sine values.  The displayed loop fixes the order of the
summation.
\begin{verbatim}
x = [pi' c];  y = [x*x]
t = 1;  u = x;  g = t;  h = u
for j = 1,...,50:
    t = [-t*y/(2*j*(2*j-1))]
    u = [-u*y/(2*j*(2*j+1))]
    g = g+t;  h = h+u
return (g,h)
\end{verbatim}
Write $g_A(m)=g(c_A(m))$ and $h_A(m)=h(c_A(m))$.

The only expanding root close to the unit circle is $\alpha$ and its
conjugate.  To treat their long tail, set
\begin{equation}\label{cert:weak-coefficients}
 a=1/\alpha,\qquad u_0=-a/p'(\alpha),\qquad
 u_1=u_0(a^3-1+1/a).
\end{equation}
Their contribution to $C_A(m)$ for $m\ge H$ is
$u_Aa^m+\overline{u_Aa^m}$.  Let $(a',u'_0)$ be the pair in
\eqref{cert:base-pairs} belonging to $z_\alpha$.  Compute
\[
 a'_0=1,\qquad a'_{n+1}=[a'a'_n]\quad(0\le n<K),\qquad
 u'_1=[u'_0(a'_3-1+[1/a'])].
\]
The subscript in $a'_n$ is a power-array index.

For a single moment use $w=u'_A$ and $k=0$.  For a paired moment
of types $A,B\in\{0,1\}$, sign $\sigma\in\{1,-1\}$, and lag
$0\le k\le K$, use
$w=u'_A+\sigma[u'_Ba'_k]$.  In either case the following calculation
returns an approximate moment.  It retains the finite factors
$-R-k\le m<H$, omits the rapidly convergent negative tail, and treats
the positive tail by summing the degree-$2$, $4$, and $6$ terms of
$-\log\cos$ through geometric series.  Their sum is the variable
\texttt{pen}; the divisors $2,12,45$ are the corresponding Taylor
denominators.  To compute $e^{-\texttt{pen}}$, the algorithm first
approximates $e^{-\texttt{pen}/32}$ and then squares five times.
This tail value is computed first, and the finite factors are then
multiplied in increasing order of $m$.  Section~\ref{cert:error}
bounds all omissions and rounding errors.
\begin{verbatim}
U = [w*a'_H]
q_0 = 1;  q_i = [q_{i-1}*U]       for i = 1,...,6
b_0 = 1;  b_i = [b_{i-1}*pi']     for i = 1,...,6
pen = 0
for d = 2,4,6:
    mom = 0
    for j = 0,...,d:
        num = [q_j*conj(q_{d-j})]
        den = 1-[a'_j*conj(a'_{d-j})]
        mom = mom+binomial(d,j)*Re([num/den])
    pen = pen+[b_d*mom/f_d]       (f_2,f_4,f_6) = (2,12,45)
x = [-pen/32];  t = 1;  value = t
for j = 1,...,40:
    t = [t*x/j];  value = value+t
repeat five times: value = [value*value]
for m = -R-k,...,H-1, in increasing order:
    v = g_A(m)                               (single)
    v = [g_A(m)*g_B(m+k)-sigma*h_A(m)*h_B(m+k)] (pair)
    value = [value*v]
return value
\end{verbatim}
Denote the outputs by $\widetilde P_A$ and
$\widetilde P_{AB}^{\sigma}(k)$, respectively.

We use only lags $30\le k\le2000$.  With $s=0$ for $\sigma=1$ and $s=1$ for
$\sigma=-1$, put $r=\widetilde P_{00}^{\sigma}(k)$ and set
\begin{equation}\label{cert:weights}
 w_{ks}=\frac{\operatorname{sgn}(r)}{W}
 \left\lfloor .68W\max(0,|r|-.00026)+\frac12\right\rfloor,
 \qquad 30\le k\le2000.
\end{equation}
Take $\operatorname{sgn}(0)=0$ and omit zero weights from subsequent
sums.  The recurrences above and this formula define the weights
without any choices or numerical tolerances.

\subsection{A uniform error bound for the moment calculation}
\label{cert:error}

We prove
\begin{equation}\label{cert:delta}
 |\widetilde P_A-P_A|<\Delta,
 \qquad
 |\widetilde P_{AB}^{\sigma}(k)-P_{AB}^{\sigma}(k)|<\Delta,
 \qquad \Delta=5\cdot10^{-6},
\end{equation}
for all types, signs, and lags used above.  The important issue is the
weakly contracting base $a=1/\alpha$: truncating its powers at $H=100$
would not be accurate.  The geometric sums in the tail calculation
retain all of these powers.

\paragraph{Coefficient and trigonometric errors.}
For $|z|<2.001$, termwise differentiation gives
\[
 |p''(z)|<4\cdot10^{13}.
\]
Hence $p'$ changes by less than $4\cdot10^{-22}$ between a center
and its root.  By~\eqref{cert:residual}, reciprocal differentiation
and rounding in~\eqref{cert:base-pairs} give base errors less than
$2/S$ and coefficient errors less than $10^{-21}$.  Direct squared
modulus comparisons give $|a_z|<1$ and $|v_z|\le1$ for every stored
pair.  The corresponding exact coefficients have modulus less than
$1.001$.  Thus each recurrence step $v\leftarrow[va_z]$ propagates
the existing error with factor at most $1$ and adds less than $4/S$.
There are fewer than $5100$ such steps in any needed entry.  Each
coefficient-times-power term consequently has error less than
$2\cdot10^{-21}$, and
\begin{equation}\label{cert:coefficient-errors}
 |c_0(m)-C_0(m)|<7\cdot10^{-20},\qquad
 |c_1(m)-C_1(m)|<2.1\cdot10^{-19}.
\end{equation}

The finite recurrence also gives the following rational comparisons:
\begin{equation}\label{cert:array-bounds}
 \begin{gathered}
 \max_{A,m}|c_A(m)|<.847,\qquad
 \max_{m\ge-20}|c_0(m)|<.018,\qquad
 \max_A|u'_A|<.0087,\\
 .8460<\max_{A,m}|c_A(m)|<.8462.
 \end{gathered}
\end{equation}
The maxima here range over the computed entries; use squared
moduli for the complex coefficients $u'_A$.  To reproduce them, perform
the coefficient recurrences in~\eqref{cert:base-pairs}, the additions
in~\eqref{cert:shift-array}, and the displayed definition of $u'_1$.
These comparisons require neither evaluation of trigonometric
functions nor a root-finding routine.

Every later trigonometric input is a sum of at most four of these
coefficients, counted with multiplicity.  Its error is less than
$10^{-18}$, and its rounded angle satisfies $|x|\le13$ and $|y|\le170$.
The angle error is less than $4\cdot10^{-18}$.  For the Taylor
recurrences at a fixed value of $y$, each term-rounding error
contributes at most $e^{14}/S$ to the final sum: the subsequent
factorial denominators bound the amplification by
$\sum_{j\ge0}170^j/(2j)!<e^{14}$.  Replacing $x^2$ by $y$ costs
at most $100e^{14}/S$, by differentiating the absolute term sums.
The omitted Taylor tails at $|x|\le13$ are less than $10^{-30}$.
Together these effects give errors below $10^{-20}$ relative to
the computed angle, and below $10^{-16}$ relative to the true
cosine and sine.  In particular every single or paired finite
factor has error less than $10^{-15}$.

There are at most $R+K+H=5100$ finite factors.  The true factors have
modulus at most $1$, and their approximations have modulus at most
$1+10^{-15}$.  Telescoping the product, including its roundings,
shows that with an initial value of modulus at most $2$ these errors
contribute less than $3\cdot10^{-11}$.

\paragraph{Terms omitted at the two ends.}
On the negative side, $|C_0(m)|\le32(.994)^{-1-m}$ for $m<0$,
by summing the contracting-root terms.  Including the three shifts
for type $1$, a single or paired coefficient at $m<-R-k$ is bounded by
\[
 200(.994)^{-4-(m+k)}.
\]
Using $1-\cos u\le u^2/2$ and telescoping products, the total omitted
product costs at most
\[
 \frac{3.142^2\,200^2}{2}
 \frac{(.994)^{5994}}{1-(.994)^2}<5\cdot10^{-8}.
\]
For example, the rational inequality $(.994)^{2996}<3\cdot10^{-8}$
already suffices for this estimate.  On the positive side the only
terms omitted from the coefficients are those from $\beta$.
Their bases are less than $.545$, and their type-$0$ and type-$1$
coefficients have modulus at most $4$.  The cosine Lipschitz bound
therefore gives a total error at most
$3.142\cdot8\cdot(.545)^{100}/(1-.545)<10^{-20}$.

\paragraph{The geometric tail and its rounding.}
Let $w_*=u_A$ in the single case and
$w_*=u_A+\sigma u_Ba^k$ in the paired case.  Put $U_*=w_*a^H$.
The exact sums of the even powers of the weak-root tail are
\begin{equation}\label{cert:geometric-moments}
 \mu_d=\sum_{n\ge0}(U_*a^n+\overline{U_*a^n})^d
 =\sum_{j=0}^d\binom dj\Re
 \frac{U_*^j\overline{U_*}^{d-j}}{1-a^j\bar a^{d-j}}
 \quad(d=2,4,6).
\end{equation}
The equality follows by the binomial theorem and absolutely
convergent geometric series.  In particular $\mu_d\ge0$.

The power-array errors satisfy $|a'_n-a^n|\le3n/S$.  Because
$|a|,|a'|>.99$, the formula for $u'_1$ gives errors below
$5\cdot10^{-21}$ for both $u'_A$, and $|u_A|\le.00871$.
It follows that the computed $U$ has error below $10^{-19}$;
its powers through degree $6$ have errors below $10^{-18}$.
In each quotient in the algorithm, the numerator error is below
$3\cdot10^{-18}$ and the denominator error is below $10^{-28}$.
The true and computed denominators have modulus greater than
$.00019$: use
\[
 |1-a^j\bar a^{d-j}|\ge1-|a|^d
 >1-(.9999)^2>.00019
\]
and the stated rounding error.  Hence each quotient error is below
$2\cdot10^{-14}$, the error in each computed $\mu_d$ is below
$2\cdot10^{-12}$, and $|\mu_d|<4\cdot10^5$ by the same denominator
bound.  The rounded powers of $\pi'$ through degree $6$ have errors
below $10^{-22}$.  Consequently \texttt{pen} approximates the
nonnegative number
\begin{equation}\label{cert:penalty}
 P=\frac{\pi^2\mu_2}{2}+\frac{\pi^4\mu_4}{12}
       +\frac{\pi^6\mu_6}{45}
\end{equation}
to within $10^{-9}$.

For $m\ge H$, the weak coefficient has angle bounded by
\[
 3.142\cdot4\cdot.00871\cdot(.9999)^m
 <.11(.9999)^m.
\]
Summing the powers of this bound gives
\[
 P\le\frac{.11^2}{2(1-.9999^2)}
      +\frac{.11^4}{12(1-.9999^4)}
      +\frac{.11^6}{45(1-.9999^6)}<31.
\]
Thus $|\texttt{pen}|<32$ and $|x|<1.1$ in the exponential
calculation.  The $40$-term exponential recurrence has Taylor
remainder less than $10^{-40}$.  Its rounding errors, bounded using
$e^{1.1}$, are negligible compared with $10^{-9}$.  Each squaring
amplifies its existing error by at most $3$, since the true
intermediate values for the computed penalty are at most
$e^{10^{-9}}$.  Including the error in~\eqref{cert:penalty}, the
returned exponential is within $4\cdot10^{-9}$ of $e^{-P}$.

\paragraph{The analytic remainder of the tail.}
For real $z$ with $|z|\le.11$ one has
\begin{equation}\label{cert:logcos-remainder}
 \left| -\log\cos z-\frac{z^2}{2}-\frac{z^4}{12}
                         -\frac{z^6}{45}\right|\le.1z^8.
\end{equation}
Here are elementary bounds proving the stated constant.  Set
$y=z^2$, $h=1-\cos z$, and
$h_0=y/2-y^2/24+y^3/720$.  Then $0\le h\le y/2$ and
$|h-h_0|\le y^4/40320$.  In
$-\log(1-h)=h+h^2/2+h^3/3+\cdots$, the terms after the third
are bounded by $h^4/(4(1-h))<.016y^4$.  Substitution of $h_0$
in the first three terms costs less than $.00004y^4$.
The terms through degree $3$ are $y/2+y^2/12+y^3/45$;
the sum of the absolute values of the remaining coefficients is
less than $.03$, so their contribution is less than $.03y^4$
for $0\le y\le.0121$.  These bounds imply~\eqref{cert:logcos-remainder}.

Applying~\eqref{cert:logcos-remainder} to all weak tail angles gives
\[
 \left|\sum_{m\ge H}-\log\cos
       \bigl(\pi(w_*a^m+\overline{w_*a^m})\bigr)-P\right|
 \le\frac{.1(.11)^8}{1-(.9999)^8}<3\cdot10^{-6}.
\]
Both penalties are nonnegative, so the same bound applies to the
difference of their exponentials.  Combining this with the rounding
error, the finite-factor error, and the two omitted ends proves
\eqref{cert:delta}; indeed the sum of the stated errors is below
$3.055\cdot10^{-6}$.

\subsection{Correlation, variance, and the absolute Gram sum}
\label{cert:moments}\label{cert:moment-proof}

For the rest of the appendix, abbreviate
$Z_i=\mathcal Z_i^{(0)}$ and $G_i=\mathcal G_i$ from
\eqref{np:stationary} and~\eqref{np:predictor}.  All expectations are
under the stationary law of the two-sided independent bits.

The preceding calculation defines exact rational weights and bounds
each moment error uniformly.  We next sum those moments to prove
the first three estimates of Proposition~\ref{np:moments}.  A final
finite-product estimate will prove the fourth.

Write $X=-.001$, $Y=.0055$, and let $\mathcal R$ be the nonzero
rows $(w_{ks},k,s)$ with $30\le k\le2000$.  Define
\[
 b_A=\widetilde P_A,\qquad
 m_{k0}=\widetilde P_{00}^{+}(k),\quad
 m_{k1}=\widetilde P_{00}^{-}(k),\qquad
 e=\widetilde P_{01}^{+}(33).
\]
The approximated correlation is
\begin{equation}\label{cert:corr-sum}
 \widetilde c=Xb_0+Ye+\sum_{(w,k,s)\in\mathcal R}wm_{ks}.
\end{equation}
For each lag $k$, let $J_{k0}$ and $J_{k1}$ be respectively the
second (minus) and first (plus) entries of the following pair:
\[
 \begin{cases}
 (\widetilde P_{10}^{+}(k-33),\widetilde P_{10}^{-}(k-33)),&k\ge33,\\
 (\widetilde P_{01}^{+}(33-k),\widetilde P_{01}^{-}(33-k)),&k<33.
 \end{cases}
\]
The approximated squared norm and absolute Gram sum are
\begin{align}
 \widetilde v={}&X^2+Y^2+2XYb_1
 +2\sum_{(w,k,s)\in\mathcal R}w(Xb_0+YJ_{ks})
 \notag\\[-2pt]
 &+\sum_{(w,k,s),(v,l,t)\in\mathcal R}
           wv\,m_{|k-l|,\,\mathbf1_{s=t}},\label{cert:var-sum}\\
 \widetilde T={}&
 \sum_{\substack{(w,k,s),(v,l,t)\in\mathcal R\\k,l\ge50}}
       |wv\,m_{|k-l|,\,\mathbf1_{s=t}}|.\label{cert:T-sum}
\end{align}
Both double sums are over ordered pairs.  Equal signs require the
minus moment in a conjugated phase pair; opposite signs require the
plus moment.  This explains the index $\mathbf1_{s=t}$, including
the diagonal cases.

The exact finite calculations give
\begin{equation}\label{cert:exact-output}
 \begin{aligned}
 |\mathcal R|&=753,\\
 |X|+|Y|+\sum_{\mathcal R}|w|&=198442946/W,\\
 |X|+|Y|+\sum_{\substack{(w,k,s)\in\mathcal R\\k<50}}|w|
     &=11816095/W,\\
 WS\widetilde c&=211815537786494774284357596751188225,\\
 W^2S\widetilde v&=143920884250294630908494564120361905955993456,\\
 W^2S\widetilde T&=105303788239450118397416374596041413014338334.
 \end{aligned}
\end{equation}
These are integer equalities: compute the arrays and products in
Section~\ref{cert:algorithm}, round the weights by~\eqref{cert:weights},
and use the unrounded rational sums
\eqref{cert:corr-sum}--\eqref{cert:T-sum}.  In particular
\[
 211810<10^9\widetilde c<211820,\quad
 143910<10^9\widetilde v<143930,\quad
 105300<10^9\widetilde T<105310.
\]
This gives $|G_i|<.199$ and $|G_i^{(N)}(x)|<.199$ by the triangle
inequality.  The total
correlation error is at most $\Delta\cdot.199$, and the squared-norm
error is at most $\Delta\cdot.199^2$.  Thus
\[
 \mathbb E\Re(Z_iG_i)>.0002108205>.000210,
 \qquad
 \mathbb E|G_i|^2<.000144119<.000145.
\]

To prepare for the last estimate, call the constant row, the type-$1$
row, and the rows with $k<50$ the \emph{small rows}; the remaining
rows are the \emph{big rows}.  A row of weight $w_v$ has phase
$U_v=\exp(\pi\mathrm i\sum_m f_v(m)e_m)$, where
$e_m=2b_{i+m}-1$.  Its sequence $f_v$ is zero for the constant,
$C_1(m+33)$ for the type-$1$ row, and
$(1-2s)C_0(m+k)$ for a row $(w,k,s)$.  Define the true absolute
Gram sum
\begin{equation}\label{cert:true-T}
 T=\sum_{v,v'\ {\rm big}}|w_vw_{v'}\mathbb E U_v\overline{U_{v'}}|.
\end{equation}
The error between this sum and~\eqref{cert:T-sum} is at most
$\Delta\cdot.199^2$, by $||x|-|y||\le|x-y|$.  Hence
\begin{equation}\label{cert:T-bound}
 T<.000116,\qquad \sqrt T<.011.
\end{equation}
The total absolute weight of the small rows is less than $.012$.

\subsection{Finite products controlling the remaining moment}
\label{cert:finite-products}

We now bound $\mathbb E(Z_i^2G_i^2)$.  The big rows have small
coefficients on the $71$ positions $-70\le m\le0$.  More precisely,
with $t=.06$, \eqref{cert:coefficient-errors} and
\eqref{cert:array-bounds} give
\begin{equation}\label{cert:tangent-bound}
 |\tan(\pi f_v(m))|\le t
 \quad(v\hbox{ big},\ -70\le m\le0).
\end{equation}
Indeed $m+k\ge-20$ for these rows, and the coefficient bound $.018$,
together with $\sin u\le u$ and $\cos u\ge1-u^2/2$, is sufficient.

For any real sequence $V$ put $c_m=|\cos(\pi V(m))|$ and
$s_m=|\sin(\pi V(m))|$.  The finite products needed below are
\begin{align}
 B_2(V)&=\prod_{m=-70}^0
  \min\left(1,\frac{(1+t^2)c_m+2ts_m}{1-t^2}\right),\notag\\
 B_1(V)&=\prod_{m=-70}^0
  \min\left(1,\frac{c_m+ts_m}{\sqrt{1-t^2}}\right),\label{cert:B-def}\\
 B_0(V)&=\prod_{m=-200}^0c_m.\notag
\end{align}
Their bounds are
\begin{equation}\label{cert:B-bounds}
 B_2(2C_0)<.011,\qquad
 \max_{f\ {\rm small}}B_1(2C_0+f)<.015,\qquad
 \max_{f,g\ {\rm small}}B_0(2C_0+f+g)<.0032.
\end{equation}
Here the maxima range over the actual phase sequences of the small
rows.  Besides the constant and the type-$1$ row, their eight
nonzero rows are as follows:
\[
\begin{array}{c|rrrrrrrr}
 k&32&36&37&39&40&43&44&47\\
 s&0&1&0&1&0&0&1&1\\
 Ww_{ks}&-536932&1265024&331701&47119&1483338&436861&458216&756904
\end{array}
\]

For completeness, the following upper-rounding calculation verifies
all of~\eqref{cert:B-bounds}.  For each of its three cases and each
indicated choice of small rows, form the approximate sequence $V'$
by replacing $C_A$ by $c_A$.  Initialize $v=1$.  For increasing $m$
in the relevant interval, compute
\begin{align*}
 b&=|g(V'(m))|+10^{-15},& y&=|h(V'(m))|+10^{-15},\\
 F_0&=b,&
 F_1&=S^{-1}+[1.002(b+.06y)],\\
 &&F_2&=S^{-1}+[1.004(1.0036b+.12y)],\\
 v&\leftarrow S^{-1}+[v\min(1,F_j)]&&\text{in case }j.
\end{align*}
The factors $1.002$ and $1.004$ bound
$(1-t^2)^{-1/2}$ and $(1-t^2)^{-1}$, respectively.  The added
$10^{-15}$ exceeds the trigonometric error proved above, and each
added $S^{-1}$ compensates for a possibly downward rounding.
Thus this calculation gives upper bounds, including for $B_0$
because its true factors are at most $1$.  The maximum final
numerators at scale $S$, in the order $B_2,B_1,B_0$, are
\begin{equation}\label{cert:B-output}
 \begin{gathered}
 10861284406435197401851134503,\\
 14731866010915969694738277129,\\
 3155364376927094122833206055.
 \end{gathered}
\end{equation}
Multiplication by $10^6/S$ bounds them respectively by $10862$,
$14732$, and $3156$, proving~\eqref{cert:B-bounds}.

\subsection{From the finite products to the anisotropy bound}
\label{cert:walsh}

The estimates just obtained concern individual coefficients and
finite products.  We finish by showing how they control the full
sum of phases without expanding every three-phase moment.
Write $G_i=G_{\rm small}+G_{\rm big}$ according to the preceding
partition, and fix any subset $J\subset\{-70,\ldots,0\}$.
Expanding the big phases in the independent signs $e_m$ on $J$ gives
\begin{equation}\label{cert:walsh-expansion}
 G_{\rm big}=\sum_{A\subset J}\mathrm i^{|A|}e_AH_A,
 \qquad e_A=\prod_{m\in A}e_m,
\end{equation}
where
\[
 H_A=\sum_{v\ {\rm big}}w_vU_v^{\rm out}
       \prod_{m\in J}\cos(\pi f_v(m))
       \prod_{m\in A}\tan(\pi f_v(m)).
\]
The phase $U_v^{\rm out}$ depends only on signs outside $J$, so
each $H_A$ is independent of the signs on $J$.

We claim that
\begin{equation}\label{cert:H-bound}
 \|H_A\|_2\le\sqrt T\,t^{|A|}(1-t^2)^{-|J|/2}.
\end{equation}
For two big rows, independence gives their full Gram entry as the
outside Gram entry times
\[
 \prod_{m\in J}\cos(\pi f_v(m))\cos(\pi f_{v'}(m))
       \bigl(1+\tan(\pi f_v(m))\tan(\pi f_{v'}(m))\bigr).
\]
By~\eqref{cert:tangent-bound}, each last factor is at least $1-t^2$.
In the expansion of $\|H_A\|_2^2$, the extra tangent products have
modulus at most $t^{2|A|}$.  Taking absolute values term by term
therefore bounds the squared norm by
$Tt^{2|A|}(1-t^2)^{-|J|}$, proving~\eqref{cert:H-bound}.

First consider $\mathbb E(Z_i^2G_{\rm big}^2)$ and set $V=2C_0$.
For a pair $A,A'\subset J$ in~\eqref{cert:walsh-expansion},
averaging the signs on $J$ contributes a factor of modulus $c_m$
where membership in $A,A'$ agrees, and a factor of modulus $s_m$
where it differs.  The outside expectation is bounded by
$\|H_A\|_2\|H_{A'}\|_2$ using Cauchy--Schwarz and the unit modulus
of the outside part of $Z_i^2$.  Summing over the four membership
choices at each site gives
\[
 |\mathbb E(Z_i^2G_{\rm big}^2)|
 \le T\prod_{m\in J}\frac{(1+t^2)c_m+2ts_m}{1-t^2}.
\]
Choose $J$ to contain precisely those positions whose displayed
factor is less than $1$.  This yields
\begin{equation}\label{cert:big-square}
 |\mathbb E(Z_i^2G_{\rm big}^2)|\le T B_2(2C_0).
\end{equation}

For a fixed small phase $U$ with sequence $f$, use the same expansion
with $V=2C_0+f$.  There is now one subset $A$; summing its two
membership choices produces $c_m+ts_m$.  The outside expectation
is bounded by $\|H_A\|_2$.  Optimizing $J$ as before gives
\[
 |\mathbb E(Z_i^2UG_{\rm big})|
 \le\sqrt T B_1(2C_0+f).
\]
For two small phases with sequences $f,g$, independence directly
bounds their full moment by $B_0(2C_0+f+g)$: discard all the other
absolute cosine factors, which are at most $1$.

Finally sum these bounds with the absolute row weights.  Using
\eqref{cert:T-bound}, \eqref{cert:B-bounds}, and the small-weight
bound $.012$, we obtain
\[
 \begin{split}
 |\mathbb E(Z_i^2G_i^2)|
 &\le T\cdot.011+2(.012)\sqrt T\cdot.015
                           +(.012)^2\cdot.0032\\
 &<.000116\cdot.011+2\cdot.012\cdot.011\cdot.015
                           +.012^2\cdot.0032\\
 &=.0000056968<.000007.
 \end{split}
\]
This proves the remaining estimate of Proposition~\ref{np:moments}.
